\begin{textsample}{POS Dim 3 – vem\_ed\_04 – Score 5.00 – vem\_ed\_04/t013.txt}  \label{ex:f3_pos_147}
PROCURA-SE Quer desenvolver um intercâmbio virtual?
Instituições de Ensino estrangeiras buscam professores nas Fatecs em diversas áreas, como: Competências interculturais Gerentes internacionais Mineração e geoprocessamento Mapa de riscos Saúde Meio Ambiente e Sustentabilidade Artes Álgebra e Cálculo História e Ciência Política \textbf{Caso} tenha interesse, escreva para cesu.pci@cps.sp.gov.br e faça parte das seguintes equipes no Teams.
Um ano de \textbf{realizações} A Coordenação de Línguas se orgulha das \textbf{realizações} neste 2020 desafiador.
A equipe dos Projetos Colaborativos Internacionais demonstrou resiliência e determinação.
No primeiro semestre, mesmo com a interrupção de atividades entre março e abril, 240 alunos, 7 professores de 8 Fatecs e 4 IES estrangeiras participaram de intercâmbios virtuais.
No segundo semestre, esses números aumentaram para 833, 35, 20 e 17, respectivamente.
A revista acadêmica CBTecLE, de Letras e Linguística, publicou 24 artigos no primeiro semestre e 11 no segundo semestre, mais um especial com o Instituto Federal de São Carlos (16 artigos).
O Congresso CBTecLE teve 10 horas de sessões online, com 12 palestrantes, 32 videopôsteres e 323 certificados emitidos.
Dedicação que se concretizou em resultados: temos muito o que comemorar.
QUEBRA-GELO Esta edição 4 traz, na seção"Quem e Quem", o depoimento de Jon Rubin, criador da iniciativa COIL (Collaborative Online International Learning) na State University of New York (SUNY).
Uma colaboração com a Universidad Tecnológica de Chile (Inacap) permitiu a \textbf{realização} de uma série de webinários sobre logística com as Fatecs Americana, Baixada Santista e Sorocaba e um PCI com a Fatec Campinas.
Na seção ' Boas práticas ', dicas para planejar PCIs, com José Carlos Barbosa Lopes, da Fatec Ipiranga.
Boa leitura!

% matched lemmas: caso, realização
\end{textsample}
